\chapter{1979年全国卷（文）}
\section{解答题}

\begin{problem}
求函数 \(y = 2x^{2} - x + 1\) 的极小值.
\end{problem}

\begin{solution}
配方得 \(y=2x^{2}-x+1=2\left(x-\frac{1}{4}\right)^{2}+\frac{7}{8}\). 因为 \(2\left(x-\frac{1}{4}\right)^{2}\geqslant0\)，所以对任意实数 \(x\)，都有 \(y\geqslant\frac{7}{8}\). 当且仅当 \(x=\frac{1}{4}\) 时等号成立，因此函数的极小值为 \(\frac{7}{8}\)，此时 \(x=\frac{1}{4}\).
\end{solution}

\begin{problem}
化简：\(\left[(1 + \sin^2\theta)^2 - \cos^4\theta\right]\left[(1 + \cos^2\theta)^2 - \sin^4\theta\right]\).
\end{problem}

\begin{solution}
由平方差公式及 \(\sin^2\theta+\cos^2\theta=1\)，原式的第一个因式为
\[
(1+\sin^2\theta)^2-\cos^4\theta=(1+\sin^2\theta-\cos^2\theta)(1+\sin^2\theta+\cos^2\theta)
\]
又 \(1+\sin^2\theta-\cos^2\theta=1+\sin^2\theta-(1-\sin^2\theta)=2\sin^2\theta\)，\(1+\sin^2\theta+\cos^2\theta=2\)，所以第一个因式等于 \(4\sin^2\theta\). 第二个因式也使用平方差公式，且有
\[
(1+\cos^2\theta)^2-\sin^4\theta=(1+\cos^2\theta-\sin^2\theta)(1+\cos^2\theta+\sin^2\theta)
\]
又 \(1+\cos^2\theta-\sin^2\theta=1+\cos^2\theta-(1-\cos^2\theta)=2\cos^2\theta\)，\(1+\cos^2\theta+\sin^2\theta=2\)，所以第二个因式等于 \(4\cos^2\theta\).
因此原式 \(=4\sin^2\theta\cdot4\cos^2\theta=16\sin^2\theta\cos^2\theta=4\sin^2 2\theta\).
\end{solution}

\begin{problem}
甲，乙二容器内都盛有酒精，甲有 \(V_{1}\text{ 公斤}\)，乙有 \(V_{2}\text{ 公斤}\). 甲中纯酒精与水（重量）之比为 \(m_{1}:n_{1}\)，乙中纯酒精与水之比为 \(m_{2}:n_{2}\). 问将二者混合后所得液体中纯酒精与水之比是多少？
\end{problem}

\begin{solution}
甲中纯酒精与水的总质量为 \(V_{1}\)，且二者的质量和对应于 \(m_{1}+n_{1}\)，所以甲中纯酒精的质量为 \(\dfrac{m_{1}V_{1}}{m_{1}+n_{1}}\)，水的质量为 \(\dfrac{n_{1}V_{1}}{m_{1}+n_{1}}\). 乙中纯酒精与水的质量和对应于 \(m_{2}+n_{2}\)，所以乙中纯酒精的质量为 \(\dfrac{m_{2}V_{2}}{m_{2}+n_{2}}\)，水的质量为 \(\dfrac{n_{2}V_{2}}{m_{2}+n_{2}}\).

混合后纯酒精的总质量为
\[
M_{\text{酒精}}=\frac{m_{1}V_{1}}{m_{1}+n_{1}}+\frac{m_{2}V_{2}}{m_{2}+n_{2}}=\frac{m_{1}V_{1}(m_{2}+n_{2})+m_{2}V_{2}(m_{1}+n_{1})}{(m_{1}+n_{1})(m_{2}+n_{2})}
\]
混合后水的总质量为
\[
M_{\text{水}}=\frac{n_{1}V_{1}}{m_{1}+n_{1}}+\frac{n_{2}V_{2}}{m_{2}+n_{2}}=\frac{n_{1}V_{1}(m_{2}+n_{2})+n_{2}V_{2}(m_{1}+n_{1})}{(m_{1}+n_{1})(m_{2}+n_{2})}
\]
两种总质量的分式具有相同的正分母，在比值中约去该分母，故混合后纯酒精与水之比为
\[
\bigl[m_{1}V_{1}(m_{2}+n_{2})+m_{2}V_{2}(m_{1}+n_{1})\bigr]:\bigl[n_{1}V_{1}(m_{2}+n_{2})+n_{2}V_{2}(m_{1}+n_{1})\bigr]
\]
\end{solution}

\begin{problem}
叙述并证明勾股定理.
\end{problem}

\begin{solution}
勾股定理：在直角三角形 \(ABC\) 中，若 \(\angle C=90^\circ\)，则斜边 \(AB\) 的平方等于两条直角边 \(AC\)、\(BC\) 的平方和，即 \(AB^2=AC^2+BC^2\).

过直角顶点 \(C\) 作斜边 \(AB\) 的垂线，垂足为 \(D\)，即 \(CD\perp AB\). 因为 \(\angle ADC=90^\circ=\angle ACB\)，且 \(A\)，\(D\)，\(B\) 共线，所以 \(\angle CAD=\angle CAB\)，从而 \(\triangle ACD\sim\triangle ABC\). 按对应边的顺序，得到 \(\dfrac{AC}{AB}=\dfrac{AD}{AC}\)，即 \(AC^2=AB\cdot AD\).

对于三角形 \(BCD\)，\(\angle BDC=90^\circ=\angle ACB\)，且 \(B\)，\(D\)，\(A\) 共线，所以 \(\angle CBD=\angle CBA\)，从而 \(\triangle BCD\sim\triangle ABC\). 因而 \(\dfrac{BC}{AB}=\dfrac{BD}{BC}\)，即 \(BC^2=AB\cdot BD\).

由于 \(D\) 在线段 \(AB\) 上，\(AD+DB=AB\). 将上面两式相加，得 \(AC^2+BC^2=AB\cdot AD+AB\cdot BD=AB(AD+BD)=AB^2\)，所以勾股定理得证.
\end{solution}

\begin{problem}
外国船只除特许外不得进入离我海岸线 \(D\text{ 里}\) 以内的区域. 设 \(A\) 及 \(B\) 是我们的观测站，\(A\) 及 \(B\) 间的距离为 \(S\text{ 里}\)，海岸线是过 \(A\)，\(B\) 的直线，一外国船在 \(P\) 点，在 \(A\) 站测得 \(\angle BAP = \alpha\)，同时在 \(B\) 站测得 \(\angle ABP = \beta\). 问 \(\alpha\) 及 \(\beta\) 满足什么简单的三角函数值不等式，就应当向此未经特许的外国船发出警告，命令退出我海域？
\end{problem}

\begin{solution}
设从船所在点 \(P\) 向海岸线 \(AB\) 作垂线，垂足为 \(C\)，并设船到海岸线的距离为 \(PC=d\). 在三角形 \(ABP\) 中，\(\angle BAP=\alpha\)，\(\angle ABP=\beta\)，故 \(\angle APB=180^\circ-(\alpha+\beta)\). 由正弦定理，\(AP=\dfrac{S\sin\beta}{\sin(\alpha+\beta)}\). 另有 \(BP=\dfrac{S\sin\alpha}{\sin(\alpha+\beta)}\).

三角形 \(ABP\) 的面积一方面为 \(\dfrac{1}{2}Sd\)，另一方面为 \(\dfrac{1}{2}AP\cdot BP\sin(\alpha+\beta)\). 代入上面的两条边长，得
\[
\frac{1}{2}Sd=\frac{1}{2}\left(\frac{S\sin\beta}{\sin(\alpha+\beta)}\right)\left(\frac{S\sin\alpha}{\sin(\alpha+\beta)}\right)\sin(\alpha+\beta)
\]
所以 \(d=\dfrac{S\sin\alpha\sin\beta}{\sin(\alpha+\beta)}\). 于是
\[
\frac{S}{d}=\frac{\sin(\alpha+\beta)}{\sin\alpha\sin\beta}=\frac{\sin\alpha\cos\beta+\cos\alpha\sin\beta}{\sin\alpha\sin\beta}=\cot\alpha+\cot\beta
\]

船进入离海岸线 \(D\text{ 里}\) 以内的区域，当且仅当 \(d\leqslant D\). 因为 \(S\)、\(d\)、\(D\) 均为正数，这个条件等价于 \(\dfrac{S}{d}\geqslant\dfrac{S}{D}\)，因此应满足的不等式为
\[
\boxed{\cot\alpha+\cot\beta\geqslant\frac{S}{D}}
\]
\end{solution}

\begin{problem}
美国的物价从 \(1939\text{ 年}\) 的 \(100\) 增加到四十年后 \(1979\text{ 年}\) 的 \(500\)，如果每年物价增长率相同，问每年增长百分之几？（注意：\(x<0.1\)，可用：\(\ln(1 + x)\approx x\)，取 \(\lg 2 = 0.3\)，\(\ln 10 = 2.3\)）
\end{problem}

\begin{solution}
设每年的物价增长率为 \(x\)，则 \(x>0\). 从 \(1939\) 年到 \(1979\) 年共经过 \(40\) 年，所以物价满足 \(100(1+x)^{40}=500\)，即 \((1+x)^{40}=5\). 两边取自然对数，得到 \(40\ln(1+x)=\ln5\).

由 \(\lg 5=\lg10-\lg2=1-0.3=0.7\)，再由换底关系 \(\ln5=\ln10\cdot\lg5\)，得 \(\ln5=2.3\times0.7=1.61\). 因为题设允许使用 \(\ln(1+x)\approx x\)，所以 \(x\approx\dfrac{1.61}{40}=0.04025\). 这个数确实小于 \(0.1\)，与近似条件相符，因此每年的物价增长率约为 \(0.04025\times100\%=4.025\%\)，按题目精度约为 \(4\%\).
\end{solution}

\begin{problem}
设 \(CEDF\) 是一个已知圆的内接矩形，过 \(D\) 作该圆的切线与 \(CE\) 的延长线相交于点 \(A\)，与 \(CF\) 的延长线相交于点 \(B\). 求证：\(\frac{BF}{AE} = \frac{BC^{3}}{AC^{3}}\).

\bitmapfigure[float=inline,width=8.5cm]{img/1979/national_liberal/q4_fig1.png}
\end{problem}

\begin{solution}
连接 \(C\)，\(D\). 因为 \(CEDF\) 是圆内接矩形，\(\angle CFD=90^\circ\). 由圆周角定理的逆定理，弦 \(CD\) 是该圆的直径，因此直线 \(CD\) 是过 \(D\) 的半径所在直线. 又因为 \(AB\) 是圆在 \(D\) 点的切线，所以 \(CD\perp AB\).

由 \(A\)，\(C\)，\(E\) 共线、\(B\)，\(C\)，\(F\) 共线，以及矩形相邻边 \(CE\perp CF\)，可得 \(\angle ACB=90^\circ\). 因而在直角三角形 \(ABC\) 中，\(CD\) 是从直角顶点 \(C\) 向斜边 \(AB\) 作出的高. 由 \(\triangle ACD\sim\triangle ABC\)，有 \(\dfrac{AC}{AB}=\dfrac{AD}{AC}\)，所以 \(AC^2=AD\cdot AB\). 由 \(\triangle BCD\sim\triangle ABC\)，有 \(\dfrac{BC}{AB}=\dfrac{BD}{BC}\)，所以 \(BC^2=BD\cdot AB\). 两式相除，得到 \(\dfrac{BD}{AD}=\dfrac{BC^2}{AC^2}\).

由点 \(A\) 的切线割线定理，\(AD^2=AE\cdot AC\)；由点 \(B\) 的切线割线定理，\(BD^2=BF\cdot BC\). 因此
\[
\frac{BF}{AE}=\frac{BD^2/BC}{AD^2/AC}=\left(\frac{BD}{AD}\right)^2\frac{AC}{BC}=\left(\frac{BC^2}{AC^2}\right)^2\frac{AC}{BC}=\frac{BC^3}{AC^3}
\]
命题得证.
\end{solution}

\begin{problem}
过原点 \(O\) 作圆 \(x^{2} + y^{2} - 2x - 4y + 4 = 0\) 的任意割线交圆于 \(P_{1}\)，\(P_{2}\) 两点，求 \(P_{1}P_{2}\) 的中点 \(P\) 的轨迹.
\end{problem}

\begin{solution}
设所作直线的一个单位方向向量为 \((\cos\varphi,\sin\varphi)\)，则直线上的点可参数表示为 \(x=t\cos\varphi\)，\(y=t\sin\varphi\). 代入圆方程，并使用 \(\cos^2\varphi+\sin^2\varphi=1\)，得到
\[
t^2-2t(\cos\varphi+2\sin\varphi)+4=0
\]
记 \(k=\cos\varphi+2\sin\varphi\)，设上式的两个根为 \(t_1\)，\(t_2\)，它们分别对应交点 \(P_1\)，\(P_2\). 由根与系数的关系，\(t_1+t_2=2k\)，所以弦中点 \(P\) 的坐标为
\[
P=\left(\frac{t_1+t_2}{2}\cos\varphi,\frac{t_1+t_2}{2}\sin\varphi\right)=(k\cos\varphi,k\sin\varphi)
\]
设 \(P=(x,y)\)，则 \(x=k\cos\varphi\)，\(y=k\sin\varphi\)，从而 \(x^2+y^2=k^2=k(\cos\varphi+2\sin\varphi)=x+2y\). 配方得
\[
\left(x-\frac{1}{2}\right)^2+(y-1)^2=\frac{5}{4}
\]

还需确定这条圆上的实际轨迹范围. 直线与原圆有两个不同交点的条件是上面关于 \(t\) 的二次方程判别式为正，即 \(4(k^2-4)>0\)，所以 \(|k|>2\). 直线的方向可以反向而不改变这条直线及弦中点，故取方向使 \(k>0\)，于是 \(k>2\). 又由柯西不等式，\(k=\cos\varphi+2\sin\varphi\leqslant\sqrt{5}\)，所以 \(2<OP=k\leqslant\sqrt{5}\).

反过来，圆 \(\left(x-\frac{1}{2}\right)^2+(y-1)^2=\frac{5}{4}\) 上满足 \(OP>2\) 的点，取 \(OP\) 所在直线作为所作直线，即可令 \((\cos\varphi,\sin\varphi)=\left(\frac{x}{OP},\frac{y}{OP}\right)\). 由该圆的方程有 \(x^2+y^2=x+2y\)，故此时 \(k=(x+2y)/OP=OP>2\)，二次方程判别式为正，确实得到两个交点. 因此上述范围没有多余点.

原圆可化为 \((x-1)^2+(y-2)^2=1\)，其圆心为 \(M(1,2)\). 轨迹圆是以 \(OM\) 为直径的圆；对其上任意点 \(P\)，有 \(OP^2+MP^2=OM^2=5\). 因而 \(OP>2\) 等价于 \(MP<1\). 所以所求轨迹为
\[
\left(x-\frac{1}{2}\right)^2+(y-1)^2=\frac{5}{4}
\]
且
\[
(x-1)^2+(y-2)^2<1
\]
即以 \(OM\) 为直径的圆位于原圆内部的一段弧，原圆与轨迹圆的两个交点（对应切线）不属于轨迹.
\end{solution}
